\chapter{Giới thiệu và quy trình phát triển}\label{ch:gioithieu}
\cref{tab:doichieu-rutgon} cho biết yêu cầu của Assignment 2 nằm ở chương nào; bản đầy đủ ở \cref{app:doichieu}.
\begin{table}[htbp]\centering
\caption{Đối chiếu rút gọn giữa Task của đề và chương báo cáo}\label{tab:doichieu-rutgon}
\begin{tabular}{@{}p{6.4cm}p{6.6cm}@{}}\toprule
\textbf{Task của đề} & \textbf{Chương} \\ \midrule
Task 1: khai thác và đặc tả yêu cầu & \cref{ch:yeucau} \\
Task 2: mô hình hóa (activity, sequence, class) & \cref{ch:mohinh} \\
Task 3: kiến trúc và triển khai (deployment) & \cref{ch:kientruc} \\
Task 4: Sprint 1, repository, MVP, CI & \cref{ch:quanly}, \cref{ch:trienkhai-se} \\
Task 5: Sprint 2, CI/CD, bàn giao & \cref{ch:trienkhai-se}, \cref{ch:kiemthu}, \cref{ch:banggiao} \\
Truy vết yêu cầu & \cref{ch:truyvet} \\ \bottomrule
\end{tabular}
\end{table}

\section{Bối cảnh, mục tiêu, phạm vi}\label{sec:se-boicanh}
Quán cà phê quy mô nhỏ cần để khách đặt món và thanh toán không dùng tiền mặt ngay trên điện thoại, tránh xếp hàng và nhầm đơn. BrewLite là ứng dụng web mobile-first cho phép xem menu, chọn size và topping, quản lý giỏ hàng, đăng nhập, đặt đơn, thanh toán mô phỏng và xem lịch sử đơn.

Assignment 2 tiếp nối Assignment 1 và đi xuyên suốt từ khai thác yêu cầu, mô hình hóa, thiết kế kiến trúc, hiện thực hai Sprint, CI/CD đến bàn giao bản production. Mã nguồn dùng chung với Assignment 1; báo cáo này tập trung vào đặc tả, thiết kế, quy trình và vận hành. Phạm vi chi tiết (trong, ngoài, không thực hiện) ở \cref{sec:se-phamvi}.

\section{Mô hình quy trình áp dụng}\label{sec:se-quytrinh}
Nhóm áp dụng Agile Scrum~\cite{scrumguide2020}: Product Backlog, Sprint Backlog, Increment; bốn sự kiện Sprint Planning, Daily Scrum, Sprint Review, Sprint Retrospective. Scrum phù hợp vì thời gian ngắn (khoảng một tuần), yêu cầu có thể điều chỉnh sau mỗi Sprint và cần sản phẩm chạy được sớm để triển khai thử. Assignment 2 chia thành hai Sprint: Sprint 1 (20--23/10) tạo nền tảng và MVP menu, Sprint 2 (24--26/10) hoàn thiện luồng end-to-end, CI/CD và bàn giao (xem \cref{sec:se-sprint}).

\section{Công nghệ và công cụ sử dụng}\label{sec:se-congnghe}
\begin{table}[htbp]\centering
\caption{Công nghệ và công cụ}\label{tab:congnghe}
\begin{tabular}{@{}p{3.2cm}p{9.6cm}@{}}\toprule
\textbf{Lớp} & \textbf{Công nghệ} \\ \midrule
Frontend & Next.js (App Router), mobile-first, manifest PWA cơ bản \\
Backend & NestJS, class-validator, JWT, bcrypt \\
Cơ sở dữ liệu & PostgreSQL 16, Prisma 7 (adapter \texttt{pg}), migration và seed \\
Kiểm thử & Vitest, Supertest (PostgreSQL thật), Playwright \\
Vận hành & Docker Compose, GitHub Actions, GHCR, Render \\
Sơ đồ & PlantUML (nguồn trong \path{docs/diagrams}) \\ \bottomrule
\end{tabular}
\end{table}
Kho mã nguồn: \url{https://github.com/tthan-prog/BrewLite}. Web production: \url{https://brewlite-web.onrender.com}. API: \url{https://brewlite-api.onrender.com/health}. Phiên bản bàn giao: tag \texttt{v0.1.2} (commit \texttt{7b1ea96}, image \texttt{sha-7b1ea96}).

\section{Cấu trúc báo cáo}\label{sec:se-cautruc}
\cref{ch:quanly} trình bày quản lý dự án; \cref{ch:yeucau} đặc tả yêu cầu; \cref{ch:mohinh} mô hình UML; \cref{ch:kientruc} kiến trúc; \cref{ch:trienkhai-se} hiện thực và CI/CD; \cref{ch:kiemthu} kiểm thử; \cref{ch:banggiao} bàn giao và nhìn lại; \cref{ch:truyvet} ma trận truy vết; \cref{ch:ketluan-se} kết luận. Phụ lục gồm bảng đối chiếu đề bài, phân công và tư liệu bổ sung.
